\section{Clementia Dei}

\enlargethispage{2em}
\thispagestyle{empty}
\begin{verse}
The soul \& heart, all mine \& all of me,\\
Floats over small body and huge abyss\\
Easily dissolved into a hot mist\\
Or not so, to rise quickly and swiftly,

Like the glimmering, vapor breath of dawn,\\
Into heaven's bourne and the breast of God,\\
As speedy and as hidden as a nod,\\
Upon the whirling wind and garden song.

The I lives on as the body unwinds.\\
What shall keep grave me sinking down to hell?\\
(So heavy and so long dead the soul that fell.)

Long time before mere Worthies' hope, I find,\\
An older good, an ancient Love, perhaps\\
The Light Who spoke to me before the Hap.
\end{verse}


\flrightit{Posted on 2023-02-23 by Logres }

\begin{center}* * *\end{center}

\begin{footnotesize}\begin{sffamily}



\texttt{Charlotte Cowell on 2012-02-11 at 05:26 said: }

That's beautiful Logres – I hear you! not easy coming to terms with the Fall is it…..thankfully the Lord is indeed merciful Cxx


\hfill

\texttt{Charlotte Cowell on 2012-02-11 at 05:54 said: }

relevant to this – and also tangentially to some of the posts about music and geometry – is an extract from MoTT Letter VI, The Lover, which I transcribed here:

\url{http://alchemical-weddings.com/alchemical-weddings/ecstasy-and-enstasy}

In particular Logres: The descent into the depths of your own soul in meditating upon the account of paradise in Genesis will render you incapable of doubt. Such is the nature of the certainty that one can have here….

….It expresses in symbolic language the first layer (first in the sense of the root of all that is human in human nature) of human psychic life, or its `beginning'. Now, knowledge of the beginning, initium in Latin, is the essence of initiation.

The above pertains to the Hermetic initiation (by which a state of enstasy is reached), achieved via spiritual touch, and this is `mirrored' by what is termed the Pythagorean initiation which has rapturous ecstasy as it's result and is attained primarily via spiritual hearing and is `essentially musical'


\hfill

\texttt{HOO on 2012-02-11 at 11:11 said: }

Impressive comment Charlotte.


\hfill

\texttt{logres on 2012-02-11 at 12:05 said: }

Agree with HOO. In fact, there might be something here as relates to Christianity's problems with traditionalism, something about beginning at the beginning. Tomberg thought that Greek and Jewish thought were both logically prior to Christianity (the Kabbalah, and the mysteries). I think he referred to them as ``Father" and ``Mother". Steiner seems to want to start with John? This helps, Charlotte, thanks.


\hfill

\texttt{Charlotte Cowell on 2012-02-11 at 13:09 said: }

the kind words are appreciated both, but I think we can all thank our unknown friend for the evergreen words of wisdom – it helps me too. 

Steiner does indeed seem to hold John in very high (highest?) esteem. He was very open in his adherence to the Rosicrucian stream, which I understand he sees as starting with John? I have also heard Steiner `accused’ – I think this is the right word! – of following John rather than Jesus, but I would not like to say that this is entirely true. What I can say I've noticed in some branches of the faith is that John is held in extreme reverence as the `originator' and this of course hearkens back to Elijah, who from a `new age gnostic' perspective is very much relevant at this point in time. 

My personal research has confirmed (to me at least) the connection between Elijah and John, but I'm not sure if this is generally accepted? All of this brings us back to Lazarus. Here is one of the many things Steiner had to say on the matter: 

``We have shown how in the course of time the being who was present in Elijah appeared again at the most important moments of human evolution on Earth—appeared again so that Christ Jesus Himself could give him the initiation he was to receive for the evolution of humankind. For the being of Elijah reappeared in Lazarus-John—who are in truth one and the same figure." 

This was his first and only indication of the connection between John the Baptist and John the Evangelist.

What else can I say except that I tend to get them `mixed up' too?!

This is what Steiner told his doctor on his deathbed:

``At the awakening of Lazarus, the spiritual Being, John the Baptist, who since his death had been the overshadowing Spirit of the disciples, penetrated from above into Lazarus as far as the Consciousness Soul; the Being of Lazarus himself, from below, intermingled with the spiritual Being of John the Baptist from above. After the awakening of Lazarus, this Being is Lazarus-John, the disciple whom Jesus loved."

I have not read teh pertinent lectures but apparently Steiner believed that after his beheading John became a kind of `group soul' for the 12 disciples, and Lazarus/John in particular. Granted the idea of Lazarus and John being one starts with Steiner, unless anyone knows otherwise?

This `sharing of spiritual essence' seems to me characteristic of Steiner's teaching, which also incorporates 2 separate Jesus children…


\hfill

\texttt{HOO on 2012-02-11 at 19:02 said: }

Speaking of poems… 

``That is why an Islamic tradition says that Adam, in the earthly Paradise, spoke in verse, that is, in rhythmic speech; this is related to that ``Syrian language" (lughah sury?niyyah) of which we spoke in our previous study on the ``science of letters," and which must be regarded as translating directly the ``solar and angelic illumination" as this manifests itself in the center of the human state. This is also why the Sacred Books are written in rhythmic language which, clearly, makes of them something quite other than the mere ``poems," in the purely profane sense, which the anti-traditional bias of the modern critics would have them to be. Moreover, in its origins poetry was by no means the vain ``literature" that it has become by a degeneration resulting from the downward march of the human cycle, and it had a truly sacred character. Traces of this can be found up to classical antiquity in the West, when poetry was still called the ``language ofthe Gods," an expression equiva lent to those we have indicated, in as much as the Gods, that is, the Devas, are, like the angels, the rep­resentation of the higher states. In Latin, verses were called carmina, a designation relating to their use in the accomplishment of rites; for the word carmen is identical to the Sanskrit karma which must be taken here in its special sense of ``ritual action"; and the poet him­self, interpreter of the ``sacred language" through which the divine Word appears, was vates […] [Guénon, The ``The Language of the Birds"]


\hfill

\texttt{Charlotte Cowell on 2012-02-12 at 05:55 said: }

yes, I read this last night, it is a very interesting text. Actually I've been enjoying everything of Guenon since I (finally!)started reading a few weeks ago – a man after my own heart in many ways, it seems. I agree with everything said about `poetry'. it's a theme that goes way back into the ancient world, from the stories of Enki to the utterances of the Pythia, Orphic hymns, epics of Homer and Virgil and so on. More recently, looking at the three interlocking circles of tradition, I always thought that the Islamic essence was very much about bringing the fire of beauty and joy to the overall equation, in no small part via the arts. It always struck me as significant that the Koran was sung and it is very relaxing to hear the call to prayer in Muslim countries. It drives some people mad when they go out there but I think it's wonderful. More specific to the Devas, there are some wonderful recitals of prayers and mantras to be heard around Diwali.

``This is she! the peerless princess, Rama's consort loved and lost,? This is she! the saintly Sita, by a cruel fortune crost,"

Hanuman thus thought and pondered: ``On her graceful form I spy,? Gems and gold by sorrowing Rama oft depicted with it sigh,

On her ears the golden pendants and the tiger's sharpened tooth,? On her arms the jewelled bracelets, tokens of unchanging truth,

On her pallid brow and bosom still the radiant jewels shine,? Rama with a sweet affection did in early days entwine!"


\hfill

\texttt{Santiago on 2023-02-26 at 00:47 said: }

This is an incredible poem and a superb delivery.


\end{sffamily}\end{footnotesize}
